% ============================================================
%  appendix/proofs.tex
%  Appendix A --- Proof of Theorem 1 (Inversion)
%  Appendix B --- Proof of Theorem 2 (Grade-Difference Group)
%  STATUS: Write W1
% ============================================================

% Note: this file is \input{} from main.tex inside \appendix.
% The \section headings live in main.tex; only proof bodies here.

% ============================================================
%  APPENDIX A --- Theorem 1: Full Proof
% ============================================================

\section{Proof of Theorem~\ref{thm:inversion} (Inversion)}
\label{app:proofs-th1}

\subsection*{Setup}

Recall: $\Phi(n,m) = 2^{n-m}$ for $(n,m)\in\mathbb{N}^2$,
$\mathrm{VACUUM} = \Phi(0,0) = 1$, and $\mathrm{Op4}$ is multiplication.

\begin{proof}[Proof of Theorem~1 (Inversion)]
We prove existence and uniqueness of $x^{-1}$ for every non-vacuum element.

\medskip
\textbf{Step 1 --- Construction.}
For $x = \Phi(n,m) = 2^{n-m}$ with $(n,m)\neq(0,0)$, define
\[
  x^{-1} \coloneqq 2^{-(n-m)} \in \mathcal{F}^*.
\]
This element exists in $\mathcal{F}^*$ since any $(a,b) \in \mathbb{N}^2$ with $a-b = -(n-m)$ indexes it.

\textbf{Step 2 --- Verification.}
\[
  \mathrm{Op4}(x, x^{-1})
  = 2^{n-m} \cdot 2^{-(n-m)}
  = 2^{(n-m)-(n-m)}
  = 2^0
  = 1
  = \Phi(0,0)
  = \mathrm{VACUUM}. \quad\checkmark
\]

\textbf{Step 3 --- Uniqueness.}
Suppose $\Phi(a,b)\cdot \Phi(n,m) = \Phi(0,0) = 1$.
Then $2^{a-b}\cdot 2^{n-m} = 2^0$, so $(a-b)+(n-m)=0$, giving
$a-b = -(n-m)$.
Hence the grades $(a,b)$ are exactly those indexing $x^{-1}$, making it unique in $\mathcal{F}^*$. $\square$

\medskip
\textbf{Worked example.}
Grade $(3,2)$: $F = 2^{3-2} = 2$.
Inverse: $F^{-1} = 2^{-(3-2)} = 2^{-1} = 0.5$.
Check: $2 \times 0.5 = 1 = \mathrm{VACUUM}$.\quad\checkmark
\end{proof}

% ============================================================
%  APPENDIX A.1 --- Proposition 1.4: Commutative Monoid (added P3.2)
% ============================================================

\subsection*{Proof of Proposition~\ref{prop:monoid} (Commutative Monoid)}
\label{app:proofs-monoid}

\begin{proof}[Proof of Proposition~\ref{prop:monoid}]
We prove that $(\Gamma, +_\Gamma)$ with componentwise addition is a
commutative monoid with identity $(0,0)$.

\medskip
\textbf{Step 1 --- Closure.}
For $(n,m),(n',m')\in\mathbb{N}^2$, $(n,m)+_\Gamma(n',m') = (n+n', m+m')$.
Since $n,n',m,m'\in\mathbb{N}$ and $\mathbb{N}$ is closed under addition,
$(n+n', m+m')\in\mathbb{N}^2 = \Gamma$. \checkmark

\medskip
\textbf{Step 2 --- Associativity.}
For $(n_1,m_1),(n_2,m_2),(n_3,m_3)\in\Gamma$:
\begin{align*}
  [(n_1,m_1)+_\Gamma(n_2,m_2)] +_\Gamma (n_3,m_3)
  &= (n_1+n_2, m_1+m_2) +_\Gamma (n_3,m_3) \\
  &= (n_1+n_2+n_3, m_1+m_2+m_3) \\
  &= (n_1,m_1) +_\Gamma [(n_2,m_2)+_\Gamma (n_3,m_3)].
\end{align*}
Follows from associativity of $+$ on $\mathbb{N}$. \checkmark

\medskip
\textbf{Step 3 --- Commutativity.}
$(n,m)+_\Gamma(n',m') = (n+n', m+m') = (n'+n, m'+m) = (n',m')+_\Gamma(n,m)$.
Follows from commutativity of $+$ on $\mathbb{N}$. \checkmark

\medskip
\textbf{Step 4 --- Identity.}
For any $(n,m)\in\Gamma$:
$(n,m)+_\Gamma(0,0) = (n+0,m+0) = (n,m)$.
So $(0,0) = \mathrm{VACUUM}$ is the identity element. \checkmark

Hence $(\Gamma, +_\Gamma)$ is a commutative monoid. $\square$
\end{proof}

% ============================================================
%  APPENDIX B --- Theorem 2: Grade-Difference Group (REWRITE)
%  Corrected result: group isomorphism, not field.
% ============================================================

\section{Proof of Theorem~\ref{thm:gdg} (Grade-Difference Group)}
\label{app:proofs-th2}

\begin{proof}[Proof of Theorem~2 (Grade-Difference Group)]

Let $\mathcal{F}^* \coloneqq \{\Phi(n,m)=2^{n-m} : (n,m)\in\mathbb{N}^2,\,
(n,m)\neq(0,0)\}$ with operation $\cdot$ (multiplication).
Define $\varphi\colon \mathcal{F}^* \to \mathbb{Z}$ by
$\varphi(\Phi(n,m)) = n - m$.

\medskip
\textbf{Homomorphism.}
$\varphi(\Phi(n,m)\cdot \Phi(n',m')) = \varphi(2^{n-m}\cdot 2^{n'-m'})
= \varphi(2^{(n+n')-(m+m')}) = (n+n')-(m+m') = (n-m)+(n'-m')
= \varphi(\Phi(n,m)) + \varphi(\Phi(n',m'))$. \checkmark

\medskip
\textbf{Bijectivity.}
\emph{Injective:} if $n_1-m_1 = n_2-m_2$ then $\varphi(F_1)=\varphi(F_2)$.
\emph{Surjective:} for any $k\in\mathbb{Z}$ choose $(n,m)=(k,0)$ if $k\geq 0$
or $(0,-k)$ if $k<0$; then $\varphi(\Phi(n,m))=k$. \checkmark

\medskip
\textbf{Identity.}
The multiplicative identity of $\mathcal{F}^*$ is $\Phi(1,1) = 2^{1-1} = 1$.
Note that VACUUM $= \Phi(0,0)$ is \emph{excluded} from $\mathcal{F}^*$, so the
group identity is $\Phi(1,1)$, not VACUUM.
$\varphi(\Phi(1,1)) = 1-1 = 0$, the additive identity of $(\mathbb{Z},+)$.
\checkmark

\medskip
\textbf{Surjectivity.}
For any $k\in\mathbb{Z}$ we exhibit a pre-image in $\mathcal{F}^*$:
\[
  \text{if } k>0:\; (k,0)\in\mathbb{N}^2\setminus\{(0,0)\},\quad
  \varphi(k,0)=k; \qquad
  \text{if } k<0:\; (0,-k)\in\mathbb{N}^2\setminus\{(0,0)\},\quad
  \varphi(0,-k)=k;
\]
\[
  \text{if } k=0:\; (1,1)\in\mathbb{N}^2\setminus\{(0,0)\},\quad
  \varphi(1,1)=0.
\]
In all three cases the chosen element lies in $\mathcal{F}^*$ and maps to $k$.
\checkmark

\medskip
Hence $\varphi\colon (\mathcal{F}^*, \cdot) \xrightarrow{\sim} (\mathbb{Z},+)$
is a group isomorphism, proving Theorem~2.
The group $G_\Delta \coloneqq (\mathcal{F}^*, \cdot)$ is the
\emph{grade-difference group}. $\square$
\end{proof}
